\documentclass[11pt]{article}
\usepackage[margin=1in]{geometry}
\usepackage{amsmath,amssymb,amsthm,mathtools}
\usepackage{hyperref}
\usepackage{enumitem}
\usepackage{booktabs}
\usepackage{xcolor}

\hypersetup{colorlinks=true, linkcolor=blue, citecolor=blue, urlcolor=blue, pdfusetitle}

\title{Entanglement-Entropy Sourcing of Gravity: Dimensional Repair,\\Covariant Completion, and Experimental Outlook}
\author{Kevin Monette}
\date{March 19, 2026}

\begin{document}
\maketitle

\begin{abstract}
We repair a dimensional inconsistency in the entanglement--gravity framework and present a refined, self-consistent phenomenological model. The original coupling prefactor $c^4/(8\pi Gk_B\ln 2)$ multiplying an entropy density in bits per cubic meter does not yield an energy density; the missing ingredient is a reference energy scale $E_*$ converting bit counts to joules. We introduce the minimal external-source effective action with coupling $\Lambda = \tilde{\kappa}E_*$, vary with respect to the metric only, and derive a dimensionally correct entanglement stress--energy tensor, exact fifth-force law, and $1/R^2$ exterior acceleration scaling for bounded sources. All experimental observables reduce to the product $\Lambda_{\mathrm{eff}} = \tilde{\kappa}E_*\alpha_{\mathrm{screen}}$; neither $\tilde{\kappa}$ nor $E_*$ is individually constrained by measurement. We identify and address a critical constraint from macroscopic superconductors, which requires the operational definition of internal entanglement entropy to exclude thermally mixed correlations. Under the Planck-scale hypothesis $E_* \approx E_P$, near-field force sensing with cryogenic SQUID-readout mechanics is the preferred experimental direction, with a predicted unsuppressed signal of $3.6\times10^{-9}$\,m/s$^2$ for 100 entangled bits at $100\,\mu$m separation.
\end{abstract}

\section{Purpose and Scope}

This document is a dimensional-consistency revision of the entanglement--gravity program. Its goal is not to select a unique microscopic origin for the coupling, but to isolate the minimal scaling assumptions required for a well-posed phenomenological theory.

Earlier drafts used a coupling prefactor proportional to $c^4/(8\pi G k_B \ln 2)$ multiplying an entanglement entropy density measured in bits per unit volume. In SI units, that expression does not yield an energy density, so the resulting stress--energy tensor and derived accelerations were not dimensionally valid. The present revision fixes that problem by introducing a reference energy scale $E_*$ and rebuilding the source term around it.

\section{Units and the Original Problem}

The Einstein equations
\begin{equation}
G_{\mu\nu} = \frac{8\pi G}{c^4} T_{\mu\nu}
\end{equation}
require the stress--energy tensor to have units of energy density,
\begin{equation}
[T_{\mu\nu}] = \mathrm{J\,m^{-3}} = \mathrm{kg\,m^{-1}\,s^{-2}}.
\end{equation}

The original framework defined the coarse-grained entanglement entropy density as
\begin{equation}
[S_{\mathrm{ent}}] = \mathrm{bit\,m^{-3}},
\end{equation}
where ``bit'' is dimensionless for unit counting. The problematic source coefficient was written as
\begin{equation}
\alpha = \tilde{\kappa}\,\frac{c^4}{8\pi G k_B \ln 2}.
\end{equation}
In SI units, this gives
\begin{equation}
[\alpha] = \mathrm{J\,m^{-1}\,K^{-1}},
\end{equation}
so that
\begin{equation}
[\alpha S_{\mathrm{ent}}] = \mathrm{J\,m^{-4}\,K^{-1}},
\end{equation}
which is not an energy density. A missing physical scale is therefore required to convert entropy density into energy density.

\section{Minimal Dimensional Fix}

\subsection{Reference energy scale}

We introduce a reference energy scale $E_*$ with units of energy per bit and define the entanglement energy density by
\begin{equation}
\rho_{\mathrm{ent}}(x) = E_*\,S_{\mathrm{ent}}(x).
\end{equation}
Then
\begin{equation}
[\rho_{\mathrm{ent}}] = [E_*][S_{\mathrm{ent}}] = \mathrm{J\,m^{-3}},
\end{equation}
as required.

An equivalent parameterization is
\begin{equation}
\rho_{\mathrm{ent}}(x) = \lambda_* k_B T_* s_{\mathrm{ent}}(x),
\end{equation}
where $s_{\mathrm{ent}}$ is entropy density in thermodynamic units, $T_*$ is an effective temperature scale, and $\lambda_*$ is dimensionless. In this white paper we use the simpler single-scale notation $E_*$.

\subsection{Revised source term}

The entanglement contribution to the stress--energy tensor is rewritten as
\begin{equation}
T^{\mathrm{ent}}_{\mu\nu}(x) = \tilde{\kappa}\,\rho_{\mathrm{ent}}(x) g_{\mu\nu} + \text{(gradient terms)},
\end{equation}
with
\begin{equation}
\rho_{\mathrm{ent}}(x) = E_* S_{\mathrm{ent}}(x).
\end{equation}
Thus all previous appearances of $\alpha S_{\mathrm{ent}}$ are replaced by $\tilde{\kappa} E_* S_{\mathrm{ent}}$.

If desired, one may later relate $E_*$ to a microscopic scale such as
\begin{equation}
E_* = \eta k_B T_*,
\end{equation}
with $\eta$ dimensionless and $T_*$ an effective Unruh or horizon temperature. The present framework leaves the origin of $E_*$ open.

\section{Revised Effective Action and Fifth-Force Equation}

\subsection{External-source action}

Let
\begin{equation}
\phi(x) \equiv S_{\mathrm{ent}}(x).
\end{equation}
We retain the external-source interpretation: $\phi$ is prescribed by the device quantum state and is not varied as an independent field. The entanglement-sector Lagrangian density is written as
\begin{equation}
\mathcal{L}_{\mathrm{ent}} = -\frac{\beta}{2} g^{\mu\nu}(\nabla_\mu\phi)(\nabla_\nu\phi) - \lambda\phi,
\end{equation}
with
\begin{equation}
\lambda = \tilde{\kappa} E_*.
\end{equation}
The full action is therefore
\begin{equation}
S[g;\phi] = \int d^4x\,\sqrt{-g}\left[\frac{c^4}{16\pi G}R - \frac{\beta}{2} g^{\mu\nu}(\nabla_\mu\phi)(\nabla_\nu\phi) - \tilde{\kappa} E_*\phi\right].
\end{equation}

Varying with respect to $g_{\mu\nu}$ yields
\begin{equation}
T^{\mathrm{ent}}_{\mu\nu} = \beta\left(\nabla_\mu\phi\,\nabla_\nu\phi - \frac{1}{2}g_{\mu\nu}(\nabla\phi)^2\right) + \tilde{\kappa} E_*\phi\, g_{\mu\nu}.
\end{equation}
The dimensions are correct provided $\beta$ carries the dimensions needed to make the gradient term an energy density. In slowly varying configurations, the gradient term is subleading.

\subsection{Matter nonconservation}

Total conservation requires
\begin{equation}
\nabla^\mu T^{\mathrm{total}}_{\mu\nu} = 0,
\end{equation}
so that
\begin{equation}
\nabla^\mu T^{\mathrm{matter}}_{\mu\nu} = -\nabla^\mu T^{\mathrm{ent}}_{\mu\nu}.
\end{equation}
Substituting the revised entanglement tensor gives
\begin{equation}
\nabla^\mu T^{\mathrm{matter}}_{\mu\nu} = -\left(\tilde{\kappa} E_* + \beta\Box\phi\right)\nabla_\nu\phi + \beta\nabla_\nu\left(\frac{1}{2}(\nabla\phi)^2\right).
\end{equation}
In the quasi-static regime with modest gradients and $|\beta\Box\phi| \ll |\tilde{\kappa}E_*|$, this reduces to
\begin{equation}
\nabla^\mu T^{\mathrm{matter}}_{\mu\nu} \approx -\tilde{\kappa}E_*\,\nabla_\nu\phi(x) = -\tilde{\kappa}E_*\,\nabla_\nu S_{\mathrm{ent}}(x).
\end{equation}
This is the corrected fifth-force law. Its units are now those of force density, as required.

\section{Newtonian Limit and Revised Observable}

\subsection{Modified Poisson equation}

In the Newtonian limit,
\begin{equation}
\nabla^2 \Phi = 4\pi G\left(\rho_{\mathrm{matter}} + \rho_{\mathrm{ent}}\right),
\end{equation}
with
\begin{equation}
\rho_{\mathrm{ent}}(\mathbf{x}) = \tilde{\kappa} E_* S_{\mathrm{ent}}(\mathbf{x}).
\end{equation}
Under spherical symmetry,
\begin{equation}
\Delta a(R) = -\frac{d\Phi_{\mathrm{ent}}}{dR} = \frac{G}{R^2}M_{\mathrm{ent}}(<R),
\end{equation}
where
\begin{equation}
M_{\mathrm{ent}}(<R) = 4\pi\int_0^R \rho_{\mathrm{ent}}(r)r^2\,dr.
\end{equation}
For bounded sources, $M_{\mathrm{ent}}(<R)$ saturates outside the source and the correction scales as $1/R^2$.

\subsection{Observable combination}

Experiments constrain the product
\begin{equation}
\Lambda_{\mathrm{eff}} \equiv \tilde{\kappa}\,E_*\,\alpha_{\mathrm{screen}},
\quad [\Lambda_{\mathrm{eff}}] = \mathrm{J/bit},
\end{equation}
not $\tilde{\kappa}$ or $E_*$ individually. In branch-comparison experiments the
useful observational combination is
\begin{equation}
\mathcal{B}
\equiv |\Lambda_{\mathrm{eff}}|\,\frac{\Delta S_{\mathrm{ent}}^{\mathrm{tot}}}{R_{\mathrm{eff}}^2},
\quad [\mathcal{B}] = \mathrm{J\,m^{-3}},
\end{equation}
where $\Delta S_{\mathrm{ent}}^{\mathrm{tot}}$ is the differential internal entanglement content between
experimental branches and $R_{\mathrm{eff}}$ is the characteristic geometric scale. A null result
at sensitivity $\Delta a_{\max}$ constrains
\begin{equation}
|\Lambda_{\mathrm{eff}}| < \frac{\Delta a_{\max}\,R_{\mathrm{eff}}^2}{G\,\Delta S_{\mathrm{ent}}^{\mathrm{tot}}}.
\end{equation}
Mapping this bound to $|\tilde{\kappa}|$ alone requires an independent assumption about
$E_*$; without such an assumption, only $|\Lambda_{\mathrm{eff}}|$ is constrained.

\section{Implications for OP8 and OP10}

\subsection{Fifth-force bounds}

The corrected acceleration outside a bounded source of radius $R_0$ is
\begin{equation}
\Delta a(R) = \frac{G\,|\Lambda_{\mathrm{eff}}|\,\Delta S_{\mathrm{ent}}^{\mathrm{tot}}}{R^2},
\quad R > R_0.
\end{equation}
A null force bound $\Delta a_{\max}$ constrains the observable product
\begin{equation}
|\Lambda_{\mathrm{eff}}|
= |\tilde{\kappa}|\,E_*\,\alpha_{\mathrm{screen}}
\lesssim \frac{\Delta a_{\max}\,R_{\mathrm{eff}}^2}{G\,\Delta S_{\mathrm{ent}}^{\mathrm{tot}}}.
\label{eq:Lambda-bound}
\end{equation}

\paragraph{Benchmark (unsuppressed signal).}
For the Planck-scale hypothesis $E_* = E_P = 1.96\times10^9$\,J and $|\tilde{\kappa}| = 1/4$, the
effective mass per bit is $m_{\mathrm{bit}} = |\tilde{\kappa}|E_P/c^2 = 5.44\,\mu\mathrm{g/bit}$.
For $\Delta S_{\mathrm{ent}} = 100$\,bits at $R = 100\,\mu\mathrm{m}$ with $\alpha_{\mathrm{screen}} = 1$:
\begin{equation}
\Delta a
= \frac{G \times 100 \times 5.44\times10^{-9}\,\mathrm{kg}}{(10^{-4}\,\mathrm{m})^2}
= 3.6\times10^{-9}\,\mathrm{m/s^2}.
\end{equation}
Cryogenic SQUID-based force sensors achieve $\sim 10^{-10}$\,m/s$^2$/\hspace{0.5pt}$\sqrt{\mathrm{Hz}}$,
placing this signal above the noise floor by a factor of $\sim\!36$ at 1\,Hz bandwidth.

\paragraph{Critical constraint.}
If $E_* \approx E_P$, macroscopic superconductors containing $\sim\!10^{20}$ Cooper pairs would
carry an effective entanglement mass of order $10^{11}$\,kg per gram of material, producing
enormous gravitational anomalies not observed. The framework requires that the operational
definition of $S_{\mathrm{int}}$ excludes uncontrolled thermal correlations (including BCS
Cooper-pair entanglement) for which no controlled bipartition protocol was applied. Verifying
this claim rigorously is an open problem; see Section~\ref{sec:open-problems}.

\subsection{RG interpretation}

With $E_*$ explicit, the natural renormalization-group ansatz becomes
\begin{equation}
\tilde{\kappa}_{\mathrm{lab}}E_{*,\mathrm{lab}} = \tilde{\kappa}_{\mathrm{UV}}E_{*,\mathrm{UV}}\left(\frac{E_{\mathrm{lab}}}{E_P}\right)^\gamma,
\end{equation}
where $E_P$ is the Planck energy and $\gamma$ is an effective scaling exponent. Experiments therefore constrain the infrared product $\tilde{\kappa}E_*$ rather than $\tilde{\kappa}$ alone.

\section{Open Problems}\label{sec:open-problems}

\subsection{Cooper-pair constraint and the scope of $S_{\mathrm{int}}$}

The single most important open problem is to determine whether the operational definition
of $S_{\mathrm{int}}$ --- the internal entanglement entropy of the device across a controlled,
pre-registered bipartition --- excludes uncontrolled thermal and BCS correlations.

If $E_* \approx E_P$ (the only assumption under which any near-term experiment is sensitive),
macroscopic superconductors such as niobium at cryogenic temperatures contain approximately
$10^{20}$ Cooper pairs per gram.  If each Cooper pair contributes a bit of $S_{\mathrm{int}}$
under any bipartition, the predicted effective entanglement mass is
\begin{equation}
M_{\mathrm{ent}}^{\mathrm{Nb}} \approx N_{\mathrm{CP}}\times\frac{|\tilde{\kappa}|E_P}{c^2}
= 10^{20}\times 5.44\times10^{-9}\,\mathrm{kg} \approx 5\times10^{11}\,\mathrm{kg/g}.
\end{equation}
This exceeds the chip's actual gravitational mass by $\sim\!5\times10^{14}$, producing enormous
local gravitational anomalies that have never been observed.

The resolution --- and it is the only available resolution --- is that BCS Cooper-pair
correlations do \emph{not} contribute to $S_{\mathrm{int}}$ as operationally defined.  The
framework defines $S_{\mathrm{int}}$ via a \emph{controlled, pre-registered bipartition protocol
applied to the prepared quantum state of the device}.  Thermally mixed Cooper pairs have not
been prepared into a pure entangled state via any such protocol; their density matrix
$\rho_{\mathrm{CP}} = \int d\xi\,p(\xi)|\xi\rangle\langle\xi|$ is a classical mixture over
BCS ground states, and the corresponding bipartite entropy $S_{\mathrm{int}}$ across any
geometric cut is zero for the pure-state component and does not accumulate as a gravitational
source under this definition.

This argument is plausible but is not yet a proof.  A rigorous demonstration requires showing
that for any density matrix $\rho$ in thermal equilibrium (i.e., $\rho \propto e^{-\beta H}$
for any $H$), the operationally defined $S_{\mathrm{int}}$ --- computed via the device's
measurement protocol and calibrated bipartition --- vanishes or is negligibly small.
This is \textbf{Open Problem~OP-new}: prove or disprove that $S_{\mathrm{int}} = 0$ for
all systems in thermal equilibrium, under the framework's operational definition.

\subsection{Status of $E_*$}

$E_*$ is a matching coefficient that cannot be determined from the framework itself.
It parameterizes the unknown UV physics connecting entanglement counting to energy.
Three physically motivated choices span many orders of magnitude:

\begin{center}\small
\begin{tabular}{@{}lll@{}}
\toprule
Choice & $E_*$ & $\Delta a$ (100 bits, 100\,$\mu$m, $\alpha_{\rm screen}=1$)\\
\midrule
Planck scale ($E_* = E_P$) & $1.96\times10^9$\,J & $3.6\times10^{-9}$\,m/s$^2$ \\
Qubit energy ($E_* = \hbar\omega$, 5\,GHz) & $3.3\times10^{-24}$\,J & $6\times10^{-42}$\,m/s$^2$ \\
Thermal ($E_* = k_BT$, 15\,mK) & $2.1\times10^{-25}$\,J & $4\times10^{-43}$\,m/s$^2$ \\
\bottomrule
\end{tabular}
\end{center}

Only the Planck-scale choice produces a signal above any conceivable near-term experimental
threshold.  The framework is therefore experimentally interesting only if $E_* \gtrsim 10^5$\,J,
which in turn requires resolving the Cooper-pair constraint via OP-new.

\subsection{OP3: Lindblad screening numerics}

A staged numerical program determines $\alpha_{\mathrm{screen}}(N,\tau)$ for realistic
NISQ platforms: exact Liouvillian evolution for $N \leq 10$; Kraus-mapped gate errors
for $N \leq 20$; MPO-TDVP tensor-network methods for $N \approx 50$.  The key diagnostic
is whether $\alpha_{\mathrm{screen}}(50, T_\phi) \gtrsim 10^{-2}$ for cluster-state initial
conditions.  Analytical predictions (confirmed for $N=10$) give $\alpha_{\mathrm{screen}}
\approx e^{-\gamma_\phi\tau}$ for cluster states (independent of $N$) vs.\ $e^{-N\gamma_\phi\tau}$
for GHZ states, strongly motivating cluster-state protocols.

\subsection{OP10: Planck-scale renormalization}

Whether $\tilde{\kappa} = -1/4$ at the Planck scale renormalizes to an observable laboratory
value depends on the scaling exponent $\gamma$ in $\tilde{\kappa}_{\mathrm{lab}}E_{*,\mathrm{lab}}
= \tilde{\kappa}_{\mathrm{UV}}E_{*,\mathrm{UV}}(E_{\mathrm{lab}}/E_P)^\gamma$.  For $\gamma = 0$
(marginal), no renormalization occurs.  For $\gamma = 4$, the laboratory coupling is suppressed
by $(E_{\mathrm{lab}}/E_P)^4 \approx 10^{-132}$.  Power counting and the classical-limit
requirement constrain $\gamma$; this problem is accessible to standard EFT techniques.

\section{Programmatic Consequences}

The dimensional repair does not alter the operational definition of $S_{\mathrm{ent}}(x)$,
the Lindblad screening program, or the computation of $\alpha_{\mathrm{screen}}(N,\tau)$.
What changes is the mapping from entanglement content to physical source strength and the
correct form of the observable combination: $\Lambda_{\mathrm{eff}} = \tilde{\kappa}E_*\alpha_{\mathrm{screen}}$
rather than $\tilde{\kappa}\alpha_{\mathrm{screen}}$ alone.

For the Planck-scale hypothesis $E_* = E_P$, near-field force sensing is the preferred
experimental direction: the acceleration signal ($3.6\times10^{-9}$\,m/s$^2$) exceeds
the sensitivity of cryogenic SQUID-based force sensors, while Ramsey-interferometry phase
shifts are suppressed by $c^2$ and lie below current thresholds by many orders of magnitude.

\section{Conclusions}

The entanglement--gravity framework is now dimensionally consistent: introducing $E_*$ repairs the units, and the resulting effective action, stress--energy tensor, matter nonconservation law, and observable combination all carry correct SI dimensions.

Three points define the current state of the framework:
\begin{enumerate}
\item \textbf{$\Lambda_{\mathrm{eff}} = \tilde{\kappa}E_*\alpha_{\mathrm{screen}}$ is the only constrained combination.}
  Experiments bound this product; separating $\tilde{\kappa}$, $E_*$, and $\alpha_{\mathrm{screen}}$
  requires additional theoretical input (OP10, OP3, and the Cooper-pair argument respectively).

\item \textbf{The Cooper-pair constraint is the critical open problem.}
  If $E_* \approx E_P$ (the only experimentally viable hypothesis), the framework requires
  that uncontrolled thermal correlations contribute zero to $S_{\mathrm{int}}$ under the
  operational definition. This must be proven, not assumed.

\item \textbf{Force sensing is the correct experimental direction.}
  For $E_* = E_P$, $|\tilde{\kappa}| = 1/4$, $\Delta S_{\mathrm{ent}} = 100$\,bits, $R = 100\,\mu$m,
  the predicted acceleration is $\Delta a = 3.6\times10^{-9}$\,m/s$^2$, within reach of
  cryogenic SQUID-based force sensors. Ramsey interferometry is suppressed by $c^2$ and
  is not the primary experimental path.
\end{enumerate}

With these three points explicit, the framework is well-posed, falsifiable, and ready for
theoretical and experimental development.

\end{document}
